\section{完整结果与复现说明}
\subsection{安全载荷与单点组批}
本附录表格保留论文评阅需要的完整服务区与架次结果。货箱逐项分配、逐箱交付、通信保障区间和任务分区的完整可机读记录见随附官方格式工作簿与 CSV。数值显示精度不改变源结果的完整精度。
\longresulttable{payloads}{15 个服务区的最大安全载荷；基准返航储备为 20\%}{payloads}
\longresulttable{q1_batches}{Q1 全部单点架次；表中前缀 Q1 省略}{q1batches}
\longresulttable{frontier}{Q2 全部候选的数值指标}{frontier}
\subsection{运输与中继完整架次表}
电池简写 A-01 表示 A-BAT-01，其余同理。所有时刻从任务起点计，以 s 为单位。实际访问顺序和货箱列表保留在工作簿，避免把过宽字符串压缩为不可读的小字号。
\longresulttable{q2_schedule}{Q2 全部运输架次；架次前缀 Q2 省略}{q2schedule}
\longresulttable{q3_schedule}{Q3 全部运输架次；架次前缀 Q3 省略}{q3schedule}
\longresulttable{relay_schedule}{中继飞行、服务窗口与能耗}{relayschedule}
\longresulttable{relay_positions}{中继悬停位置与返航 SOC}{relayposition}
\subsection{计算入口与证据文件}
将官方附件原样解压至程序的数据目录后，依次运行以下命令可生成正式结果、诊断和图表。所有路径在项目根目录下解析，Windows 与 Linux 使用同一 Python 入口；论文需要含中文支持的 XeLaTeX。
\begin{verbatim}
python run_all.py
python -m pytest -q
python scripts/analyze_paper.py
python scripts/setup_paper_fonts.py
python scripts/build_paper.py
\end{verbatim}
正式结果验证文件记录官方输入的 SHA256 摘要。图形契约记录每图的核心主张、数据路径、脚本、单位与不确定性口径；所有主图同时提供 PDF、SVG 和 300 dpi PNG。新增采样诊断不覆盖正式结果，且每次纸面数字更新都应重新生成图、表和工作簿。

\section{人工智能辅助与第三方来源说明}
本工作使用 ChatGPT/Codex（OpenAI，2026-09-23）辅助代码审计、图形设计和文字草拟。提示要求以官方规则为准、不编造数据与文献；输出经数据回放、边界测试、加密通信检查和逐页审查。本次精确模型版本与公开发布日期在执行记录中不可得，未沿用无法核验的历史声明；参赛队应从实际使用界面补全官方要求的信息，并独立核验和理解全部内容。AI 建议不作为公式或实验数据来源。

科研图形组织参考 TAO-QKV 的 Icarus-Figures（\url{https://github.com/TAO-QKV/Icarus-Figures}，访问日期 2026-09-23）：采用一图一主张、突出主面板、展示约束阈值、灰度冗余编码和三格式导出。全部无人机场景图由本项目真实数据重新绘制，没有粘贴外部示例图片，也未使用生成式图像。
